% Môn: Vật lý
% Lớp: 12
% Chương: Chương II. Khí lí tưởng
% Bài: L12C2 Bài 12. Áp suất khí theo mô hình động học phân tử · Giải thích áp suất khí bằng va chạm phân tử
% ===== Câu 10 =====
% ID: L12C2B12-61-DS
% Mức: VD
\dangbt{Giải thích áp suất khí bằng va chạm phân tử}
\begin{ex}
% ID: L12C2B12-61-DS
% Mức: VD
Trong dạng “Chuyển động Brown và bằng chứng về chuyển động phân tử”, Giải thích vì sao chất khí gây áp suất lên thành bình.
\choiceTF

\loigiai{
Mỗi va chạm làm phân tử đổi động lượng và truyền xung lượng cho thành; tổng hợp rất nhiều va chạm tạo lực trung bình trên một đơn vị diện tích, tức áp suất.
}
\end{ex}

% ===== Câu 20 =====
% ID: L12C2B12-62-TN
% Mức: TH
\begin{ex}
% ID: L12C2B12-62-TN
% Mức: TH
Trong dạng “Chuyển động Brown và bằng chứng về chuyển động phân tử”, Áp suất khí lên thành bình sinh ra chủ yếu do
\choice
{\True phân tử va chạm thành bình}
{trọng lượng bình}
{lực hút của Trái Đất}
{thành bình nóng lên}
\loigiai{
Đáp án A. Xung lượng truyền trong các va chạm phân tử với thành bình tạo áp suất.
}
\end{ex}

% ===== Câu 34 =====
% ID: L12C2B12-63-DS
% Mức: VD
\begin{ex}
% ID: L12C2B12-63-DS
% Mức: VD
Trong dạng “Các luận điểm cơ bản của mô hình động học phân tử chất khí”, Giải thích vì sao chất khí gây áp suất lên thành bình.
\choiceTF

\loigiai{
Mỗi va chạm làm phân tử đổi động lượng và truyền xung lượng cho thành; tổng hợp rất nhiều va chạm tạo lực trung bình trên một đơn vị diện tích, tức áp suất.
}
\end{ex}

% ===== Câu 58 =====
% ID: L12C2B12-64-DS
% Mức: VD
\begin{ex}
% ID: L12C2B12-64-DS
% Mức: VD
Trong dạng “Các thông số trạng thái của một lượng khí”, Giải thích vì sao chất khí gây áp suất lên thành bình.
\choiceTF

\loigiai{
Mỗi va chạm làm phân tử đổi động lượng và truyền xung lượng cho thành; tổng hợp rất nhiều va chạm tạo lực trung bình trên một đơn vị diện tích, tức áp suất.
}
\end{ex}

% ===== Câu 60 =====
% ID: L12C2B12-65-TN
% Mức: TH
\begin{ex}
% ID: L12C2B12-65-TN
% Mức: TH
Trong dạng “Các luận điểm cơ bản của mô hình động học phân tử chất khí”, Áp suất khí lên thành bình sinh ra chủ yếu do
\choice
{\True phân tử va chạm thành bình}
{trọng lượng bình}
{lực hút của Trái Đất}
{thành bình nóng lên}
\loigiai{
Đáp án A. Xung lượng truyền trong các va chạm phân tử với thành bình tạo áp suất.
}
\end{ex}

% ===== Câu 61 =====
% ID: L12C2B12-66-DS
% Mức: TH
\begin{ex}
% ID: L12C2B12-66-DS
% Mức: TH
Xét các phát biểu sau trong dạng “Giải thích áp suất khí bằng va chạm phân tử”.
\choiceTF
{\True Theo mô hình động học phân tử, các phân tử khí — phát biểu đúng là: chuyển động hỗn loạn không ngừng.}
{Theo mô hình động học phân tử, các phân tử khí — phát biểu đứng yên tại vị trí cân bằng là đúng.}
{\True Kích thước phân tử so với khoảng cách trung bình giữa chúng trong khí lí tưởng được coi là — phát biểu đúng là: rất nhỏ.}
{Kích thước phân tử so với khoảng cách trung bình giữa chúng trong khí lí tưởng được coi là — phát biểu bằng nhau là đúng.}
\loigiai{
\begin{itemchoice}
\item (Đúng) Theo mô hình động học phân tử, các phân tử khí — phát biểu đúng là: chuyển động hỗn loạn không ngừng. (Vì): Phân tử khí luôn chuyển động nhiệt hỗn loạn không ngừng. b) Sai: nội dung không phù hợp. c) Đúng: Trong mô hình khí lí tưởng, kích thước phân tử được bỏ qua so với khoảng cách giữa chúng. d) Sai: nội dung không phù hợp
\item (Sai) Theo mô hình động học phân tử, các phân tử khí — phát biểu đứng yên tại vị trí cân bằng là đúng.
\item (Đúng) Kích thước phân tử so với khoảng cách trung bình giữa chúng trong khí lí tưởng được coi là — phát biểu đúng là: rất nhỏ.
\item (Sai) Kích thước phân tử so với khoảng cách trung bình giữa chúng trong khí lí tưởng được coi là — phát biểu bằng nhau là đúng.
\end{itemchoice}
}
\end{ex}

% ===== Câu 65 =====
% ID: L12C2B12-67-DS
% Mức: TH
\begin{ex}
% ID: L12C2B12-67-DS
% Mức: TH
Xét các phát biểu sau trong dạng “Giải thích áp suất khí bằng va chạm phân tử”.
\choiceTF
{\True Lực tương tác giữa các phân tử khí lí tưởng được coi là — phát biểu đúng là: không đáng kể trừ khi va chạm.}
{Lực tương tác giữa các phân tử khí lí tưởng được coi là — phát biểu luôn rất lớn là đúng.}
{\True Chuyển động Brown là chuyển động — phát biểu đúng là: hỗn loạn của hạt nhỏ lơ lửng do va chạm phân tử.}
{Chuyển động Brown là chuyển động — phát biểu rơi tự do của hạt bụi là đúng.}
\loigiai{
\begin{itemchoice}
\item (Đúng) Lực tương tác giữa các phân tử khí lí tưởng được coi là — phát biểu đúng là: không đáng kể trừ khi va chạm. (Vì): Ngoài thời gian va chạm rất ngắn, tương tác phân tử được bỏ qua. b) Sai: nội dung không phù hợp. c) Đúng: Các phân tử môi trường va chạm không cân bằng làm hạt nhỏ chuyển động Brown. d) Sai: nội dung không phù hợp
\item (Sai) Lực tương tác giữa các phân tử khí lí tưởng được coi là — phát biểu luôn rất lớn là đúng.
\item (Đúng) Chuyển động Brown là chuyển động — phát biểu đúng là: hỗn loạn của hạt nhỏ lơ lửng do va chạm phân tử.
\item (Sai) Chuyển động Brown là chuyển động — phát biểu rơi tự do của hạt bụi là đúng.
\end{itemchoice}
}
\end{ex}

% ===== Câu 69 =====
% ID: L12C2B12-68-DS
% Mức: VD
\begin{ex}
% ID: L12C2B12-68-DS
% Mức: VD
Xét các phát biểu sau trong dạng “Giải thích áp suất khí bằng va chạm phân tử”.
\choiceTF
{\True Áp suất khí lên thành bình sinh ra chủ yếu do — phát biểu đúng là: phân tử va chạm thành bình.}
{Áp suất khí lên thành bình sinh ra chủ yếu do — phát biểu trọng lượng bình là đúng.}
{\True Khi nhiệt độ khí tăng, chuyển động nhiệt của phân tử nhìn chung — phát biểu đúng là: nhanh hơn.}
{Khi nhiệt độ khí tăng, chuyển động nhiệt của phân tử nhìn chung — phát biểu không đổi là đúng.}
\loigiai{
\begin{itemchoice}
\item (Đúng) Áp suất khí lên thành bình sinh ra chủ yếu do — phát biểu đúng là: phân tử va chạm thành bình. (Vì): Xung lượng truyền trong các va chạm phân tử với thành bình tạo áp suất. b) Sai: nội dung không phù hợp. c) Đúng: Nhiệt độ tăng làm động năng tịnh tiến trung bình tăng. d) Sai: nội dung không phù hợp
\item (Sai) Áp suất khí lên thành bình sinh ra chủ yếu do — phát biểu trọng lượng bình là đúng.
\item (Đúng) Khi nhiệt độ khí tăng, chuyển động nhiệt của phân tử nhìn chung — phát biểu đúng là: nhanh hơn.
\item (Sai) Khi nhiệt độ khí tăng, chuyển động nhiệt của phân tử nhìn chung — phát biểu không đổi là đúng.
\end{itemchoice}
}
\end{ex}

% ===== Câu 73 =====
% ID: L12C2B12-69-DS
% Mức: VD
\begin{ex}
% ID: L12C2B12-69-DS
% Mức: VD
Xét các phát biểu sau trong dạng “Giải thích áp suất khí bằng va chạm phân tử”.
\choiceTF
{\True Ba thông số trạng thái cơ bản của một lượng khí là — phát biểu đúng là: áp suất, thể tích, nhiệt độ.}
{Ba thông số trạng thái cơ bản của một lượng khí là — phát biểu khối lượng, trọng lượng, nhiệt lượng là đúng.}
{\True Va chạm giữa các phân tử khí lí tưởng được xem là — phát biểu đúng là: đàn hồi.}
{Va chạm giữa các phân tử khí lí tưởng được xem là — phát biểu không bảo toàn động lượng là đúng.}
\loigiai{
\begin{itemchoice}
\item (Đúng) Ba thông số trạng thái cơ bản của một lượng khí là — phát biểu đúng là: áp suất, thể tích, nhiệt độ. (Vì): Trạng thái lượng khí xác định bởi p, $V$, T. b) Sai: nội dung không phù hợp. c) Đúng: Mô hình khí lí tưởng coi các va chạm là đàn hồi. d) Sai: nội dung không phù hợp
\item (Sai) Ba thông số trạng thái cơ bản của một lượng khí là — phát biểu khối lượng, trọng lượng, nhiệt lượng là đúng.
\item (Đúng) Va chạm giữa các phân tử khí lí tưởng được xem là — phát biểu đúng là: đàn hồi.
\item (Sai) Va chạm giữa các phân tử khí lí tưởng được xem là — phát biểu không bảo toàn động lượng là đúng.
\end{itemchoice}
}
\end{ex}

% ===== Câu 77 =====
% ID: L12C2B12-70-DS
% Mức: VDC
\begin{ex}
% ID: L12C2B12-70-DS
% Mức: VDC
Xét các phát biểu sau trong dạng “Giải thích áp suất khí bằng va chạm phân tử”.
\choiceTF
{\True Khi nén khí ở nhiệt độ không đổi, mật độ phân tử — phát biểu đúng là: tăng.}
{Khi nén khí ở nhiệt độ không đổi, mật độ phân tử — phát biểu giảm là đúng.}
{\True Hiện tượng khuếch tán chứng tỏ các phân tử — phát biểu đúng là: chuyển động không ngừng.}
{Hiện tượng khuếch tán chứng tỏ các phân tử — phát biểu đứng yên là đúng.}
\loigiai{
\begin{itemchoice}
\item (Đúng) Khi nén khí ở nhiệt độ không đổi, mật độ phân tử — phát biểu đúng là: tăng. (Vì): Cùng số phân tử trong thể tích nhỏ hơn nên mật độ phân tử tăng. b) Sai: nội dung không phù hợp. c) Đúng: Khuếch tán là hệ quả trực tiếp của chuyển động nhiệt hỗn loạn. d) Sai: nội dung không phù hợp
\item (Sai) Khi nén khí ở nhiệt độ không đổi, mật độ phân tử — phát biểu giảm là đúng.
\item (Đúng) Hiện tượng khuếch tán chứng tỏ các phân tử — phát biểu đúng là: chuyển động không ngừng.
\item (Sai) Hiện tượng khuếch tán chứng tỏ các phân tử — phát biểu đứng yên là đúng.
\end{itemchoice}
}
\end{ex}

% ===== Câu 82 =====
% ID: L12C2B12-71-DS
% Mức: VD
\begin{ex}
% ID: L12C2B12-71-DS
% Mức: VD
Trong dạng “Giải thích áp suất khí bằng va chạm phân tử”, Giải thích vì sao chất khí gây áp suất lên thành bình.
\choiceTF

\loigiai{
Mỗi va chạm làm phân tử đổi động lượng và truyền xung lượng cho thành; tổng hợp rất nhiều va chạm tạo lực trung bình trên một đơn vị diện tích, tức áp suất.
}
\end{ex}

% ===== Câu 100 =====
% ID: L12C2B12-72-TN
% Mức: TH
\begin{ex}
% ID: L12C2B12-72-TN
% Mức: TH
Trong dạng “Các thông số trạng thái của một lượng khí”, Áp suất khí lên thành bình sinh ra chủ yếu do
\choice
{\True phân tử va chạm thành bình}
{trọng lượng bình}
{lực hút của Trái Đất}
{thành bình nóng lên}
\loigiai{
Đáp án A. Xung lượng truyền trong các va chạm phân tử với thành bình tạo áp suất.
}
\end{ex}

% ===== Câu 106 =====
% ID: L12C2B12-73-DS
% Mức: VD
\begin{ex}
% ID: L12C2B12-73-DS
% Mức: VD
Trong dạng “Vận dụng mô hình động học phân tử vào thực tế”, Giải thích vì sao chất khí gây áp suất lên thành bình.
\choiceTF

\loigiai{
Mỗi va chạm làm phân tử đổi động lượng và truyền xung lượng cho thành; tổng hợp rất nhiều va chạm tạo lực trung bình trên một đơn vị diện tích, tức áp suất.
}
\end{ex}

% ===== Câu 127 =====
% ID: L12C2B12-74-DS
% Mức: VD
\begin{ex}
% ID: L12C2B12-74-DS
% Mức: VD
Trong dạng “Đặc điểm cấu tạo và chuyển động phân tử ở thể khí”, Giải thích vì sao chất khí gây áp suất lên thành bình.
\choiceTF

\loigiai{
Số phân tử khí không đổi nên mật độ phân tử $n = \frac{N}{V}$ tỉ lệ nghịch với thể tích $V$.
Tỉ số mật độ phân tử sau và trước khi nén là: $\frac{n_2}{n_1} = \frac{V_1}{V_2} = \frac{4}{2} = 2$.
Mật độ phân tử tăng 2 lần.

Đáp số: 2
}
\end{ex}

% ===== Câu 135 =====
% ID: L12C2B12-75-TN
% Mức: TH
\begin{ex}
% ID: L12C2B12-75-TN
% Mức: TH
Trong dạng “Giải thích áp suất khí bằng va chạm phân tử”, Áp suất khí lên thành bình sinh ra chủ yếu do
\choice
{\True phân tử va chạm thành bình}
{trọng lượng bình}
{lực hút của Trái Đất}
{thành bình nóng lên}
\loigiai{
Thể tích của khối khí được tính theo công thức $V = \frac{N}{n_0}$.
Thay số: $V = \frac{1,2 \times 10^{24}}{3 \times 10^{25}} = 0,04 \, \mathrm{m^3}$.
Đổi đơn vị: $0,04 \, \mathrm{m^3} = 40 \, \mathrm{lít}$.

Đáp số: 40 lít
}
\end{ex}

% ===== Câu 147 =====
% ID: L12C2B12-76-TN
% Mức: TH
\begin{ex}
% ID: L12C2B12-76-TN
% Mức: TH
Trong dạng “Vận dụng mô hình động học phân tử vào thực tế”, Áp suất khí lên thành bình sinh ra chủ yếu do
\choice
{\True phân tử va chạm thành bình}
{trọng lượng bình}
{lực hút của Trái Đất}
{thành bình nóng lên}
\loigiai{
Đáp án A. Xung lượng truyền trong các va chạm phân tử với thành bình tạo áp suất.
}
\end{ex}

% ===== Câu 157 =====
% ID: L12C2B12-77-TN
% Mức: TH
\begin{ex}
% ID: L12C2B12-77-TN
% Mức: TH
Trong dạng “Đặc điểm cấu tạo và chuyển động phân tử ở thể khí”, Áp suất khí lên thành bình sinh ra chủ yếu do
\choice
{\True phân tử va chạm thành bình}
{trọng lượng bình}
{lực hút của Trái Đất}
{thành bình nóng lên}
\loigiai{
Đáp án A. Xung lượng truyền trong các va chạm phân tử với thành bình tạo áp suất.
}
\end{ex}
